%% kapa: exam=1 round=37 year=2026 period=1 form=A subject=civil_law subject_ko=민법 seq=33 number=33
%% source: Q-Net 감정평가사 기출문제 (cst00309) artlSeq=5258924 file=raw/1008/5258924/2026년 제37회 감정평가사 1차 1교시.pdf
%% conversion: kordoc 4.24.51 -> tools/gen1.py (auto); review=proofread-sonnet; flags=none
\begin{question}{33}
지역권에 관한 설명으로 옳은 것은? (다툼이 있으면 판례에 따름)
\begin{choices}
  \item 지역권은 요역지와 분리하여 양도하거나 다른 권리의 목적으로 할 수 있다.
  \item 승역지 소유자는 지역권의 행사를 방해하지 않는 범위 내에서 지역권자가 지역권의 행사를 위하여 승역지에 설치한 공작물을 사용할 수 있다.
  \item 요역지 소유자가 승역지 소유자와 승역지상에 도로를 개설하여 영구히 사용하기로 약정하였다면, 존속기간을 정하지 않았으므로 지역권 설정에 관한 합의가 있었다고 볼 수 없다.
  \item 요역지 소유자는 통로의 개설 없이 승역지를 통행한 사실만으로도 지역권을 시효취득할 수 있다.
  \item 지역권자는 승역지의 점유침탈이 있는 경우, 지역권에 기하여 승역지 반환청구권을 행사할 수 있다.
\end{choices}
\end{question}
\answer{2}
